\textbf{Method.}
We formulate dialect translation as a recursive, structure-driven instantiation problem over a query graph $\mathcal{G}$. Given an abstract query tree, the system constructs the final dialect-compliant SQL query in a bottom-up manner, where leaf segments are resolved first and progressively composed into higher-level structures. This dependency-driven design ensures that nested subqueries are always grounded in already validated fragments, thereby reducing error propagation and alleviating context complexity during generation.

To guarantee that the generated queries remain semantically consistent with the target dialect during recursive construction, the system employs a rule-guided instantiation process as formalized in Algorithm~\ref{alg:translation}. The algorithm first performs a topological sort on the structure graph $\mathcal{G}$ to determine a bottom-up execution sequence (Line~\ref{line:alg2:sort}), ensuring that nested subqueries are resolved before their parents. For each segment, the algorithm recursively substitutes instantiated child fragments into the current structure (Lines~\ref{line:alg2:substitute_start}--\ref{line:alg2:substitute_end}).

Each abstract token $\tau$ is instantiated using rules retrieved from the Hybrid Knowledge Base (Line~\ref{line:alg2:retrieve}), which integrates a structural pattern store and a vector-indexed repository of dialect evidence. If the retrieval confidence falls below a threshold $\delta$, a rule generator synthesizes a new transformation rule from the AST context and few-shot examples (Line~\ref{line:alg2:generate}), which is then cached back into the knowledge base for future reuse. The token is then translated into target-compliant syntax using the selected rule (Line~\ref{line:alg2:translate}).

To ensure syntactic correctness under semantic constraints, each generated fragment is validated by the target database parser. If a parsing error occurs, an iterative repair loop is triggered, where the fragment is corrected using both the error message and structural context while preserving semantic intent (Line~\ref{line:alg2:fix}). This mechanism avoids degenerate fixes such as clause removal and enforces faithful preservation of query semantics.

Finally, validated fragments are progressively cached and assembled in a bottom-up manner to produce the final output query $Q_{target}$ (Lines~\ref{line:alg2:cache}--\ref{line:alg2:assemble}).

\textbf{Figure Interpretation.}
Figure~\ref{fig:agent_translation} illustrates the full execution pipeline. The left panel shows the abstract query graph $\mathcal{G}$ with hierarchical dependencies. The center panel corresponds to Algorithm~\ref{alg:translation}, where each node is instantiated in a bottom-up order through rule retrieval, adaptive rule generation, and iterative repair. The right panel presents the final dialect-compliant SQL query, demonstrating how abstract structural representations are progressively grounded into executable expressions.